% Generated draft: requires mathematical and source review.
\documentclass[11pt]{article}
\usepackage[margin=25mm]{geometry}
\usepackage{fontspec}
\setmainfont{Latin Modern Roman}
\usepackage{amsmath,amssymb,mathtools,unicode-math}
\setmathfont{Latin Modern Math}
\usepackage{enumitem,needspace,xurl,hyperref}
\setlength{\parindent}{0pt}
\setlength{\parskip}{5pt}
\setlength{\emergencystretch}{3em}
\newcommand{\N}{\mathbb{N}}
\newcommand{\Z}{\mathbb{Z}}
\newcommand{\Q}{\mathbb{Q}}
\newcommand{\R}{\mathbb{R}}
\newcommand{\C}{\mathbb{C}}
\newcommand{\F}{\mathbb{F}}
\newcommand{\A}{\mathbb{A}}
\newcommand{\PP}{\mathbb{P}}
\newcommand{\E}{\mathbb{E}}
\newcommand{\eps}{\varepsilon}
\newcommand{\dd}{\,\mathrm{d}}
\begin{document}
\section*{UnsolvedMath ID: 30001845}
\textbf{Source problem code:} catalogue entry
\subsection{Definitions and mathematical statement}
The supplied record refers to rational parameters \(A,B,C\) and a Nahm series \(F_{A,B,C}(q)\), but supplies neither the series formula nor the types, dimensions, and admissibility conditions of the parameters. It also leaves the meaning of modularity unspecified.

Let \(\mathcal D\) denote the intended admissible class of rational parameter triples, and let \(\operatorname{Mod}(F_{A,B,C})\) denote the intended modularity property. These are placeholders for definitions missing from the supplied record. The requested problem is to give a necessary-and-sufficient condition \(\mathcal K(A,B,C)\) satisfying
\[
\forall (A,B,C)\in\mathcal D,\qquad
\operatorname{Mod}(F_{A,B,C})\ \Longleftrightarrow\ \mathcal K(A,B,C),
\]
and to formulate the corresponding repair of Nahm's conjectural criterion.

The display expresses the requested logical form only. A precise mathematical conjecture requires specification of \(\mathcal D\), the definition and domain of \(F_{A,B,C}\), the modularity convention, and the original conjectural criterion being repaired. None of these missing choices is fixed here.
\subsection{Short English statement}
Find an exact criterion for modularity of the Nahm series with rational parameters \(A,B,C\), and explain how Nahm's conjectural criterion should be corrected. The supplied record omits definitions and hypotheses needed to make this question precise.
\subsection{Sources}
\begin{itemize}
\item Supplied local UnsolvedMath registry record, permanent ID \(30001845\), title: \emph{Correct Modularity Criteria for Nahm Series}; source problem code: \texttt{catalogue entry}. Collection: Oberwolfach Reports Open Problems. The record gives proposed year 2011 and identifies no proposer.
\item Supplied report citation: \emph{Explicit Methods in Number Theory (2012)}, Oberwolfach Reports, pp.~1971--2044, DOI \texttt{10.4171/owr/2011/35}. Supplied source URL and literature-source link, labeled ``original OWR report'': \url{https://doi.org/10.4171/owr/2011/35}. The background describes these reports as documenting mathematical workshops and their open-problem discussions.
\item Literature assessment dated \texttt{2026-08-21}: the supplied curation labels the problem open, states that the report calls the correct Nahm-type formulation an interesting open question, and asserts that later work has not supplied a general resolution of this extracted formulation. This is supplied dated curation, not an independently verified literature assessment; time-sensitive status and cited claims require independent verification.
\item The registry labels the statement \texttt{corrected\_verified} and reports that its mathematical content was checked against the official report, with notation, context, and question form normalized without changing the source-level claim. This is a supplied verification claim; the report has not been independently consulted for this draft.
\end{itemize}
\end{document}
